% Generated draft: requires mathematical and source review.
\documentclass[11pt]{article}
\usepackage[margin=25mm]{geometry}
\usepackage{fontspec}
\setmainfont{Latin Modern Roman}
\usepackage{amsmath,amssymb,mathtools,unicode-math}
\setmathfont{Latin Modern Math}
\usepackage{enumitem,needspace,xurl,hyperref}
\setlength{\parindent}{0pt}
\setlength{\parskip}{5pt}
\setlength{\emergencystretch}{3em}
\newcommand{\N}{\mathbb{N}}
\newcommand{\Z}{\mathbb{Z}}
\newcommand{\Q}{\mathbb{Q}}
\newcommand{\R}{\mathbb{R}}
\newcommand{\C}{\mathbb{C}}
\newcommand{\F}{\mathbb{F}}
\newcommand{\A}{\mathbb{A}}
\newcommand{\PP}{\mathbb{P}}
\newcommand{\E}{\mathbb{E}}
\newcommand{\eps}{\varepsilon}
\newcommand{\dd}{\,\mathrm{d}}
\begin{document}
\section*{UnsolvedMath ID: 30004507}
\textbf{Source problem code:} catalogue entry
\subsection{Definitions and mathematical statement}
The supplied record asks for subconvexity at the central point \(s=\tfrac12\) for \(L(\operatorname{sym}^{2}u_j,s)\) as the spectral parameter \(t_j\) grows. Here \(\operatorname{sym}^{2}\) denotes the symmetric-square lift. The record does not specify the family \((u_j)_j\), its underlying group, level, character, or the normalization relating \(t_j\) to its spectral eigenvalue. Consequently, a fully specified unconditional formulation cannot be recovered from this record alone.

Under the conventional interpretation that \((u_j)_j\) is an intended family of Hecke--Maass cusp forms with Laplace eigenvalues \(\tfrac14+t_j^2\), with all auxiliary arithmetic data fixed, let \(L(\operatorname{sym}^{2}u_j,s)\) denote the standard finite symmetric-square \(L\)-function in the normalization with central point \(s=\tfrac12\). Let \(Q_j\) be its analytic conductor at that point. A precise conditional formulation of the requested subconvexity bound is: there exists \(\delta>0\) such that, for every \(\varepsilon>0\), there is a constant \(C_{\varepsilon}>0\), depending on the fixed family data but independent of \(j\), for which
\[
\left|L\!\left(\operatorname{sym}^{2}u_j,\tfrac12\right)\right|
\le C_{\varepsilon}Q_j^{\,1/4-\delta+\varepsilon}
\qquad\text{for every }j.
\]
If the intended fixed-level spectral family has \(Q_j\asymp(1+|t_j|)^2\), this is equivalently a request for some \(\eta>0\) such that, for every \(\varepsilon>0\),
\[
\left|L\!\left(\operatorname{sym}^{2}u_j,\tfrac12\right)\right|
\ll_{\varepsilon,\,\mathrm{fixed\ data}}
(1+|t_j|)^{1/2-\eta+\varepsilon}
\qquad\text{for every }j.
\]
These family assumptions and conductor specialization are explicitly conditional interpretations, not hypotheses supplied by the record. The target concerns growth in the form's spectral parameter \(t_j\), with the evaluation point held at \(s=\tfrac12\).
\subsection{Short English statement}
Obtain a uniform power saving over the convexity bound for central values of symmetric-square \(L\)-functions as the forms' spectral parameters grow. The intended family and its arithmetic data are missing, so the displayed formulation is conditional on specifying them. The supplied dated curation labels the problem partially solved, reporting off-central or mixed-range results and strong moment estimates while identifying the central spectral-aspect bound as outstanding; this assessment has not been independently verified.
\subsection{Sources}
\begin{itemize}
\item Supplied local registry record in \texttt{data/unsolvedmath/problems.json}, permanent UnsolvedMath ID \(30004507\), source problem code \texttt{catalogue entry}, title: \emph{Subconvexity for Symmetric-Square L-Functions}. The record belongs to the Oberwolfach Reports Open Problems collection; no individual proposer is supplied.
\item Supplied report citation: \emph{Automorphic Forms and Arithmetic} (2021), Oberwolfach Reports, DOI \texttt{10.4171/owr/2020/26}, pp.~1317--1383. \url{https://doi.org/10.4171/owr/2020/26}. The registry records proposed year 2020. Its background describes Oberwolfach Reports as documenting mathematical workshops and their open-problem discussions.
\item Supplied literature reference: \emph{Moments and hybrid subconvexity for symmetric-square L-functions}. \url{https://arxiv.org/abs/2009.08419}.
\item Supplied literature reference: \emph{Subconvexity for symmetric-square L-functions off-centre}. \url{https://arxiv.org/abs/2302.07236}.
\item Literature assessment date, exactly as supplied: \texttt{2026-08-22}. This is supplied dated curation derived from the OWR record, not an independently verified literature assessment. It labels the problem partially solved, reports subconvexity in off-central or mixed ranges and strong moments, and identifies central-point subconvexity in the spectral parameter as the outstanding target. The source explicitly calls for independent verification of time-sensitive status and cited claims.
\item The supplied statement-status label is \texttt{corrected\_verified}; its accompanying explanation says mathematical content was preserved while notation, typography, omitted definitions, or duplicate clauses were repaired. This records the registry's editorial assessment and does not establish independent mathematical verification of this draft.
\end{itemize}
\end{document}
